\clearpage\section{Caracterización del problema (Justificación)}
El creciente interés en las energías renovables, el almacenamiento energético y la electrificación de sectores estratégicos ha impulsado la necesidad de desarrollar sistemas de conversión de potencia más eficientes, precisos y confiables. En este contexto, los convertidores DC-DC bidireccionales juegan un papel fundamental, ya que permiten gestionar el flujo de energía en ambos sentidos entre distintas fuentes y cargas de corriente continua. Estos convertidores son clave en aplicaciones como sistemas de baterías (BESS), vehículos eléctricos, sistemas fotovoltaicos y microrredes DC, donde la precisión y robustez en el control son requisitos esenciales.

Las microrredes DC presentan ventajas significativas en términos de eficiencia y compatibilidad con fuentes renovables, pero su correcto funcionamiento depende de la implementación de sistemas de control sofisticados para los convertidores asociados. Particularmente, la topología de puente medio bidireccional en configuración interleaved ha demostrado potencial para mejorar la calidad del rizado de corriente, la distribución térmica y la respuesta dinámica del sistema. Sin embargo, el diseño e implementación del lazo de control para esta configuración requiere una cuidadosa coordinación de fases y una validación rigurosa.

En este sentido, aunque existen estudios teóricos y simulaciones, la validación experimental mediante entornos Hardware-in-the-Loop (HIL) sigue siendo un área con amplio potencial de desarrollo. El uso de HIL permite realizar pruebas en tiempo real de los algoritmos de control, combinando simulaciones precisas con hardware físico, lo que facilita la identificación y corrección de fallas antes de la implementación final.

Por lo tanto, esta tesis se enfoca en la propuesta, diseño y validación de un lazo de control para un convertidor DC-DC bidireccional en configuración interleaved, utilizando una plataforma HIL. Esta aproximación permitirá evaluar integralmente el desempeño del sistema bajo condiciones realistas, contribuyendo al avance en el control electrónico de convertidores y fortaleciendo su aplicabilidad en microrredes DC y otros sistemas donde la conversión bidireccional es crítica.